\documentclass[11pt]{article}
\usepackage[a4paper,margin=29mm]{geometry}
\usepackage[T1]{fontenc}
\usepackage[utf8]{inputenc}
\usepackage{lmodern,amsmath,amssymb,amsthm,mathtools,microtype}
\usepackage[colorlinks=true,linkcolor=blue,citecolor=blue,urlcolor=blue]{hyperref}
\hypersetup{pdftitle={Vanishing center-of-mass diffusion for the three-dimensional FENE dumbbell model},pdfauthor={Kaiwen Jin},pdfsubject={Kinetic FENE fluids and the zero spatial diffusion limit}}
\numberwithin{equation}{section}
\newtheorem{theorem}{Theorem}[section]
\newtheorem{proposition}[theorem]{Proposition}
\newtheorem{lemma}[theorem]{Lemma}
\newtheorem{remark}[theorem]{Remark}
\newcommand{\T}{\mathbb T^3}
\newcommand{\R}{\mathbb R}
\newcommand{\E}{\mathbb E}
\newcommand{\dd}{\,\mathrm d}
\newcommand{\Div}{\operatorname{div}}
\newcommand{\tr}{\operatorname{tr}}
\newcommand{\Ent}{\operatorname{Ent}_M}
\newcommand{\grad}{\nabla}
\newcommand{\Id}{\mathrm{Id}}
\newcommand{\supp}{\operatorname{supp}}
\newcommand{\Var}{\operatorname{Var}}
\newcommand{\Hone}{\mathcal H_1}
\newcommand{\Htwo}{\mathcal H_2}
\newcommand{\Dq}{\mathcal D_q}
\newcommand{\Dy}{\mathcal D_y}
\newcommand{\weak}{\rightharpoonup}
\newcommand{\pair}[1]{\langle #1\rangle}
\allowdisplaybreaks[2]
\title{Vanishing center-of-mass diffusion for the three-dimensional FENE dumbbell model}
\author{Kaiwen Jin\\\small\href{mailto:yui.ui.siki@gmail.com}{\texttt{yui.ui.siki@gmail.com}}}
\date{15 October 2026}
\begin{document}
\maketitle
\begin{abstract}
We study the zero center-of-mass diffusion limit for the full kinetic FENE
dumbbell model on the three-dimensional torus. For every solvent viscosity
$\nu>0$, FENE exponent $k>2$, smooth divergence-free initial velocity, and
smooth normalized initial density with bounded positive Maxwellian quotient,
we select global positive-diffusion weak solutions converging to a global
weak solution of the identical nondiffusive system. Velocity converges strongly
in $L^2$ and polymer density and stress converge strongly in $L^1$ on finite
time intervals. All transport and configuration fluxes have their original
constitutive forms in the limit. The proof conditions on a spatial Brownian
motion common to every polymer configuration. Nonlinear configuration
entropies are transformed before expectation, giving only a linear spatial
diffusion term. Conditional stress cutoffs, a transported limiting normalizer,
and a joint covariance estimate then propagate a zero entropy defect without
assuming strong convergence of the velocity gradient.
\end{abstract}

\section*{Introduction}
In a kinetic dumbbell model, the fluid transports a probability distribution
of polymer configurations, and the corresponding elastic stress feeds back
into the fluid momentum equation. Finitely extensible nonlinear elasticity
confines the configuration to a ball by a potential singular at its boundary.
Center-of-mass diffusion describes translational Brownian motion of the
polymer molecules. Removing this term should preserve both the distribution
and its coupling to the full velocity gradient. For weak solutions at large
data, this requires a compactness assertion beyond separate existence for the
diffusive and nondiffusive equations.

Barrett and S\"uli \cite{BS} established large-data weak existence with positive
center-of-mass diffusion. Masmoudi \cite{Masmoudi} established weak existence
without that diffusion by propagating a transformed-density defect and using
stress-driven Stokes dissipation. The distinction between those existence
results and passage from one model to the other is explicitly described in
\cite[Section~1]{DS}. The result below addresses that passage for the full
noncorotational model on $\T\times B_1$, with arbitrary prescribed data in
the stated smooth bounded-positive class. No smallness condition is imposed.

The elementary estimates already give strong velocity compactness and identify
the singular stress using a configuration boundary tail. They do not identify
the product of a weakly convergent velocity gradient with a weakly convergent
configuration moment. The nondiffusive entropy-defect argument provides a
way to close this correlation, but applying its nonlinear entropy quotients
directly to a spatially diffusive density creates spatial Hessian terms whose
sign is unavailable.

We instead represent each positive-diffusion density by conditioning on one
spatial Brownian motion shared by all labels and configurations. Given this
motion, each spatial characteristic carries an ordinary configuration
Fokker--Planck equation. Its entropy and logarithmic moments can be computed
before expectation. The expected transformed equation then has just a linear
spatial Laplacian, which disappears against fixed tests as its coefficient
tends to zero. A fixed-velocity comparison identifies this representation with
the chosen coupled density.

The remaining argument keeps this conditional structure when removing entropy
cutoffs. Radial positivity preserves the monotone stress squares, and the
deterministic velocity gradient gives a common direction to the cutoff drags.
These facts identify the stress-work covariance through a local Stokes energy
comparison. A joint law of the gradient, conditional square root, and
configuration derivative retains the positive Fisher moments. An explicit
two-state inequality supplies their averaged coercivity. Finally, a normalizer
transported by the limiting fluid makes the entropy defect bounded and gives a
characteristic Gronwall inequality. This is an extension of the defect method
in \cite{Masmoudi}; the conditional construction and its averaged cutoff and
coercivity arguments are proved here rather than attributed to that work.

Section~\ref{sec:theorem} states the precise result. Sections~\ref{sec:prelim}
and \ref{sec:lift} provide the uniform estimates and identify the conditional
representation. Sections~\ref{sec:envelope}--\ref{sec:law} establish the common
normalizer, stress-work identity, and configuration coercivity. The normalized
defect closes in Section~\ref{sec:defect}, and Section~\ref{sec:final} passes
the original system and its initial-time entropy inequality to the limit.
\section{The system and the compactness theorem}\label{sec:theorem}
Let $B=B_1(0)\subset\R^3$ and let $\T$ have its usual periodic measure. Fix
$\nu>0$ and $k>2$, and set
\begin{equation}\label{eq:weight}
 U(q)=-k\log(1-|q|^2),\qquad
 M(q)=Z^{-1}(1-|q|^2)^k,\qquad Z=\int_B(1-|q|^2)^k\dd q.
\end{equation}
We consider the full velocity-gradient FENE system
\begin{align}
 \partial_tu+u\cdot\grad_xu-\nu\Delta_xu+\grad_xp
   &=\Div_x\tau(\psi),\qquad\Div_xu=0,\label{eq:NS}\\
 \tau(\psi)&=\int_B K(q)\psi\dd q,\qquad
 K(q)=q\otimes\grad_qU(q).\label{eq:stress}
\end{align}
\begin{equation}\label{eq:FP}
\begin{aligned}
 \partial_t\psi+u\cdot\grad_x\psi
   &=\varepsilon\Delta_x\psi\\
   &\quad+\Div_q\{M\grad_q(\psi/M)-(\grad_xu)q\psi\}.
\end{aligned}
\end{equation}
The configuration flux has zero normal component at $\partial B$. Its weak
meaning is the variational form of \eqref{eq:FP} for smooth tests on the closed
ball, with no configuration boundary term. Thus no value of $\psi/M$ on the
sphere is prescribed. We use $(\grad_xu)_{ij}=\partial_{x_j}u_i$; radial symmetry
makes $K$ symmetric.

Write $L^p_M=L^p(B,M\dd q)$ and
$H^1_M=\{h:\int_BM(h^2+|\grad_qh|^2)<\infty\}$. For a normalized density define
\begin{equation}\label{eq:energy}
\begin{aligned}
 \Ent(\psi)&=\int_{\T\times B}\psi\log(\psi/M),\qquad
 \mathcal E(u,\psi)=\tfrac12\|u\|_2^2+\Ent(\psi),\\
 I_q(\psi)&=\int_{\T\times B}M|\grad_q\sqrt{\psi/M}|^2.
\end{aligned}
\end{equation}
The spatial Fisher information $I_x$ is defined by replacing $\grad_q$ by
$\grad_x$. The usual initial-time energy inequality is
\begin{equation}\label{eq:entropy}
 \mathcal E(u(t),\psi(t))+\nu\int_0^t\|\grad_xu\|_2^2\dd s
 +4\int_0^t I_q(\psi)\dd s+4\varepsilon\int_0^t I_x(\psi)\dd s
 \leq \mathcal E(u_0,\psi_0).
\end{equation}
All assertions about this inequality at a positive time are understood for
almost every time, with the usual weakly continuous representatives.

\begin{theorem}\label{thm:main}
Let $u_0$ be a smooth divergence-free periodic velocity, and let
$\psi_0=Mg_0$, where $g_0$ is smooth on $\T\times B$, bounded above and below
by positive constants, and $\int_BMg_0(x,q)\dd q=1$ for every $x$.
For every $\nu>0$ and $k>2$ one can select global normalized nonnegative
finite-entropy weak solutions $(u_\varepsilon,\psi_\varepsilon)$ of
\eqref{eq:NS}--\eqref{eq:FP}, for $\varepsilon>0$, and a sequence
$\varepsilon_j\downarrow0$ such that, on every finite interval $(0,T)$,
\begin{align}
 u_{\varepsilon_j}&\longrightarrow u
       &&\text{strongly in }L^2((0,T)\times\T),\label{eq:uconv}\\
 \psi_{\varepsilon_j}&\longrightarrow\psi
       &&\text{strongly in }L^1((0,T)\times\T\times B),\label{eq:psiconv}\\
 \tau(\psi_{\varepsilon_j})&\longrightarrow\tau(\psi)
       &&\text{strongly in }L^1((0,T)\times\T).\label{eq:tauconv}
\end{align}
The pair $(u,\psi)$ is a global weak solution of the same full system with
$\varepsilon=0$. Its spatial transport, configuration drag, and configuration
diffusion fluxes are, respectively, $u\psi$, $(\grad_xu)q\psi$, and
$M\grad_q(\psi/M)$, and the corresponding positive-diffusion fluxes converge
to them in distributions. The zero configuration-flux condition and the
prescribed initial data are retained. The limiting pair satisfies
\eqref{eq:entropy} with $\varepsilon=0$ and
\[
 u\in L^\infty_{\mathrm{loc}}([0,\infty);L^2(\T))
       \cap L^2_{\mathrm{loc}}([0,\infty);H^1(\T)),\qquad
 \psi\in C_{\mathrm{weak}}([0,\infty);L^1(\T\times B)).
\]
The density has a strong $L^1$ initial trace. No smallness assumption or strong
convergence assumption on the full velocity gradient is imposed.
\end{theorem}

Only smoothness in the open configuration ball is required of $g_0$; its
derivatives need not extend to $\partial B$. If a chosen positive-diffusion
family has the energy inequality starting at almost every positive time, that
additional inequality passes to the limit. We do not include this additional
selection property in the unconditional assertion of Theorem~\ref{thm:main}.

\section{Estimates before the configuration-drag limit}\label{sec:prelim}
The fixed-positive-diffusion construction and several configuration estimates
will be used below. We spell out their spatial domain and coefficient choices,
since existence at $\varepsilon=0$ alone does not give the limit in this paper.

\subsection{Positive diffusion and periodic geometry}\label{subsec:existence}
Barrett and S\"uli \cite[Theorem~6.1]{BS} construct weak solutions at fixed
$\varepsilon>0$ by entropy truncation and time discretization. In their
one-spring notation take the scalar potential
$V(z)=-k\log(1-2z)$, $0\leq z<1/2$, so that
$V'(z)=2k/(1-2z)$. With $d(q)=1-|q|$, the Maxwellian is comparable to $d^k$ and
$k\leq dV'\leq2k$. Their potential assumptions therefore hold with exponent
$k>2$. Choose configuration diffusion coefficient one and stress coefficient
one. Their isotropic stress subtraction is $\Id$ for number density one; its
spatial divergence is zero.

The spatial boundary condition in \cite{BS} is not periodic. Here is the
adaptation of that fixed-$\varepsilon$ construction. Replace its velocity
space by the closure of mean-zero solenoidal trigonometric polynomials and its
density space by $H^1(\T\times B,M)$. The velocity projectors preserve
solenoidality, periodic integration by parts cancels transport, and mean-zero
Poincar\'e replaces the zero-boundary Poincar\'e inequality. The spatial compact
embeddings follow from periodic charts; the configuration form domain and its
compact embedding are unchanged. The finite-cutoff variational existence,
positivity, and mass tests use precisely these properties. Terms that vanished
on the spatial boundary now cancel between opposite faces. The configuration
tests and the potential inequalities remain the same.

For completeness, use the cutoff entropy with second derivative
$1/\min(s,L)$. Its configuration-drag term cancels the full Kramers stress
work. At zero body force no forcing estimate consumes viscous dissipation.
When $L>\|g_0\|_\infty$, the initial cutoff entropy equals the true initial
entropy. The coupled approximation thus retains \eqref{eq:entropy} with its
displayed constants. At fixed $\varepsilon$, the compactness and cutoff
passages in \cite{BS} use the same weighted configuration and spatial compact
embeddings on these periodic spaces. Strong $L^1$ density convergence gives
strong $L^2$ convergence of bounded configuration moments by normalization;
these moments pair with the weak $L^2$ velocity gradients. The Hardy estimate
below removes stress collars. Lower semicontinuity gives \eqref{eq:entropy}.
Globally extended approximations and nested extraction on integer intervals
give a global solution. For a nonzero initial mean $c$, translate $x$ by $ct$
and add $c$ to the mean-zero velocity. Both equations are preserved, and the
constant mean kinetic energy cancels from \eqref{eq:entropy}. This is a
periodic adaptation of a fixed-positive-diffusion construction, not an imported
vanishing-diffusion assertion.

\subsection{Configuration boundary control and velocity compactness}
We use the quadratic Hardy estimate and the weighted Sobolev gain from
\cite[Lemma~3.2(ii), Proposition~3.7]{Masmoudi}:
\begin{align}
 \int_B\frac{Mh^2}{d^2}&\leq C_k\int_BM(|\grad_qh|^2+h^2),\label{eq:Hardy}\\
 \|h\|_{L^{p_*}_M}&\leq C_k\|h\|_{H^1_M}
       \quad\text{for some }p_*>2.\label{eq:Sobolev}
\end{align}
The first holds for $k>1$. Fix a finite interval and put
$f_\varepsilon=\sqrt{\psi_\varepsilon/M}$ and
$i_\varepsilon(t,x)=\int_BM|\grad_qf_\varepsilon|^2$.
Normalization and \eqref{eq:Hardy} imply
\begin{equation}\label{eq:tails}
 \int_B\frac{\psi_\varepsilon}{d^2}\leq C(1+i_\varepsilon),\qquad
 \int_{d<r}\frac{\psi_\varepsilon}{d}\leq Cr(1+i_\varepsilon).
\end{equation}
Integration by parts, justified by these collars, gives
\begin{equation}\label{eq:tauFisher}
 \tau(\psi_\varepsilon)=\Id+
 2\int_BM q\otimes(f_\varepsilon\grad_qf_\varepsilon),\qquad
 |\tau(\psi_\varepsilon)-\Id|\leq2\sqrt{i_\varepsilon}.
\end{equation}
In particular the stress is bounded in $L^2_{t,x}$. The equilibrium term
$\Id$ in this formula is essential.

The energy inequality and momentum equation bound
$\partial_tu_\varepsilon$ in $L^{4/3}_tH^{-1}_x$, after solenoidal projection.
Indeed the three-dimensional energy interpolation controls $u\otimes u$,
and viscosity and stress use their $L^2$ bounds. Project onto finitely many
periodic Laplace modes. Every coefficient has bounded $L^{4/3}$ time
derivative and is uniformly H\"older continuous with exponent $1/4$.
A finite-dimensional subsequence converges uniformly in time. The integrated
squared tail beyond spatial frequency $L$ is at most $C/L^2$ by the gradient
bound and Bessel's inequality. A diagonal in $L$ proves strong $L^2_{t,x}$
compactness. We henceforth write
\begin{equation}\label{eq:weakstart}
 \varepsilon_n\downarrow0,\quad u_n\to u\text{ in }L^2_{t,x},\quad
 A_n:=\grad_xu_n\weak A:=\grad_xu\text{ in }L^2_{t,x},\quad
 \psi_n\weak\psi\text{ in }L^1_{t,x,q}.
\end{equation}
The last extraction follows from entropy uniform integrability on the finite
domain. Normalization passes by bounded $(t,x)$ tests. Since $|K|\leq C/d$,
\eqref{eq:tails} identifies the weak stress by bounded kernel truncation.

The spatial transport already passes: for bounded smooth $\phi$,
\[
 \left|\int(u_n-u)\psi_n\phi\right|\leq
 \|\phi\|_\infty\|u_n-u\|_{L^1_{t,x}}.
\]
In the term containing $u(\psi_n-\psi)$, truncate $u$, use weak $L^1$
convergence, then remove the truncation using unit configuration mass.
Similarly $\varepsilon_n\Delta_x\psi_n$ tends to zero against fixed tests at
order $\varepsilon_n$. What remains is the product of $A_n$ with bounded
configuration moments of $\psi_n$. Neither factor is yet strongly compact.

\section{Conditioning on the common spatial Brownian motion}\label{sec:lift}
The nonlinear quantities used to remove the remaining correlation will be
applied to a conditional density. Throughout this section the velocity $v$ is
prescribed, deterministic, divergence-free, and belongs to $L^2_tH^1_x$;
$A=\grad_xv$ and $\varepsilon>0$ are fixed.

\subsection{The fibre equation and its measurable flow representation}
Let $W$ be one standard three-dimensional Brownian motion, common to every
spatial label and every configuration. Put $Z=\sqrt{2\varepsilon}W$ and let
$Y$ be the regular Lagrangian flow of $v(t,y+Z(t))$. Then $X=Y+Z$ satisfies
\begin{equation}\label{eq:SDE}
 \dd X=v(t,X)\dd t+\sqrt{2\varepsilon}\dd W,
 \qquad X(0,b)=b.
\end{equation}
The flow preserves volume for almost every path. The deterministic Sobolev
flow facts needed here, including stability and the characteristic meaning of
bounded transport inequalities with $L^1$ source, are
\cite[Appendix~A]{Masmoudi}. Translation preserves all their spatial Sobolev
bounds. Smooth approximation gives the adapted version of this flow. The
related scalar stochastic representation on the torus is described in
\cite[Section~3.3]{BCC}; the configuration fibres below are an additional
part of the present construction.

For any trace-free matrix history $\mathsf B\in L^2(0,T)$, the weighted variational
problem is
\begin{equation}\label{eq:fibre}
 \pair{\partial_tr,h}+\int_BM\grad_qr\cdot\grad_qh
   =\int_BM\mathsf B(t)qr\cdot\grad_qh,
 \qquad r(0)=r_0\in L^2_M.
\end{equation}
A dense Galerkin basis of the form domain, including constants, gives
\begin{equation}\label{eq:fibreenergy}
 \tfrac12\frac{\dd}{\dd t}\|r\|_{L^2_M}^2
 +\tfrac12\|\grad_qr\|_{L^2_M}^2
 \leq\tfrac12|\mathsf B|^2\|r\|_{L^2_M}^2.
\end{equation}
Weak extraction yields
$r\in L^\infty_tL^2_M\cap L^2_tH^1_M$ and
$\partial_tr\in L^2_t(H^1_M)'$. Testing with the negative part proves
positivity; the constant test conserves mass. The same estimate for a
difference proves uniqueness and time regularization gives a continuous
$L^2_M$ representative. Galerkin measurability and uniqueness give measurable,
progressive solutions for measurable progressive histories.

Use $\mathsf B(t)=A(t,X(t,b))$ and $r_0=g_0(b,\cdot)$. Volume preservation and Fubini
give $\mathsf B\in L^2_t$ on almost every label and path. Define
\begin{equation}\label{eq:rho}
 \rho^\omega(t,x,q)=M(q)r^\omega(t,X_t^{-1}(x),q).
\end{equation}
Each density has unit configuration mass. Its labelwise $L^2_M$ bound may
grow exponentially with the individual history; we do not claim a global
$L^2$ bound in the labels or paths. The uniform global estimates below use
entropy instead. It\^o's formula against a smooth test of $X$, with the fibre
equation, shows that $\E\rho$ solves the prescribed-velocity version of
\eqref{eq:FP}.

\subsection{Conditional entropy and logarithmic moments}\label{subsec:log}
Applying the configuration entropy calculation to \eqref{eq:fibre}, then
integrating the labels, gives for almost every path
\begin{equation}\label{eq:conditionalentropy}
 \Ent(\rho^\omega(t))+2\int_0^t I_q(\rho^\omega)\dd s
 \leq\Ent(\psi_0)+C\int_0^t\int_\T|A|^2.
\end{equation}
The constant is deterministic: volume preservation makes the last integral
the same on every path. In deriving this estimate use \eqref{eq:tauFisher}
and absorb $2|A|\sqrt{i_q}$ into Fisher dissipation. No conditional fluid
equation is used.

Fix $a=e^4$ and introduce, for each conditional profile,
\begin{equation}\label{eq:transforms}
\begin{aligned}
 s=\rho/M+a,\quad f=\sqrt s,\quad
 G&=\Theta(s)=\sqrt{s\log s},\\
 \Hone&=\int_BMs\log s,\\
 \Htwo&=\left(\int_BMs[(\log s)^2-2\log s+2]\right)^{1/2}.
\end{aligned}
\end{equation}
The shift keeps all scalar derivatives in a fixed positive range. Along a
label $\Htwo$ is absolutely continuous with derivative
\begin{equation}\label{eq:H2source}
 F=-\frac4{\Htwo}\int_BM\log s|\grad_qf|^2
 +\frac1{\Htwo}\int_BM(\mathsf B:K)
       \left[s\log s-s-\frac a2\log^2s\right].
\end{equation}
Indeed the derivative of $s[(\log s)^2-2\log s+2]$ is $\log^2s$ and its
second derivative is $2\log s/s$. Integration by parts gives
\eqref{eq:H2source}. The unshifted drag part is alternatively
$2\Htwo^{-1}\int M\mathsf Bqf\log s\grad_qf$. Since
$\Hone\leq C\Htwo$, it is bounded by
$C|\mathsf B|(\int M\log s|\grad f|^2/\Htwo)^{1/2}$.
The shifted part is bounded by $C|\mathsf B|/\Htwo$, since
\[
 \int_B\frac{M\log^2s}{d}\leq
 \left(\int_B\frac M{d^2}\right)^{1/2}
 \left(\int_BM\log^4s\right)^{1/2}\leq C_{k,a}.
\]
Here $\int_BMs=a+1$, $k>1$, and $\log^4s\leq C_as$.
The gradient of $G$ is
$(\sqrt{\log s}+1/\sqrt{\log s})\grad f$, whose square is comparable to
$\log s|\grad f|^2$. Young's inequality, the positive lower bound on
$\Htwo$, and \eqref{eq:H2source} now give
\begin{equation}\label{eq:logbounds}
 \sup_t\int_\T\Htwo^\omega+
 \int_0^T\int_\T\frac{\int_BM|\grad_qG^\omega|^2}{\Htwo^\omega}
 +\int_0^T\int_\T|F^\omega|\leq C.
\end{equation}
All these calculations follow first from convex truncations with derivative
$\log^2(\min(s,R))$ and then from monotone convergence and Fatou. Thus they
hold in the fibre class without boundary regularity of the initial datum.

Interpolating \eqref{eq:Sobolev} with the fixed $L^2_M$ mass of $f$ gives a
full time-space-configuration exponent $r_0=4-4/p_*>2$ for $f$.
Every slightly smaller exponent above two also controls $G$.
Moreover $\Htwo^2\leq C\int_BM f^{2+\zeta}$ for sufficiently small
$\zeta>0$, so $\Htwo$ is bounded in $L^2_{t,x}$, uniformly in the path.
The moments $\Hone,\Htwo$ and $G^2$ are uniformly integrable on the joint
probability-time-space domain. These are integrability estimates, not spatial
compactness assertions.

\subsection{Identification with the coupled density}\label{subsec:identification}
We next justify that this conditional construction represents the actual
positive-diffusion coupled solution, rather than another solution of its
prescribed-velocity equation.

\begin{lemma}\label{lem:unique}
Fix $\varepsilon>0$, $k>1$, and a divergence-free $v\in L^2_tH^1_x$.
Two normalized nonnegative solutions of its linear Fokker--Planck equation
coincide if both have finite configuration and $\varepsilon$-weighted spatial
Fisher information and the same strong $L^1$ initial trace.
\end{lemma}
\begin{proof}
For densities $p_i$, put $b=(v,Aq-\grad U)$ and
$D_\varepsilon=\operatorname{diag}(\varepsilon\Id,\Id)$.
Hardy and normalization give
\[
 \sqrt{p_i}\in L^2_tH^1_{x,q},\qquad
 \int |b|^2p_i\leq C\left(\int|v|^2+\int|A|^2+
                              \int p_i/d^2\right)<\infty.
\]
For $\delta>0$, set $s_i=p_i+\delta M$ and
\[
 J(s_1,s_2)=\frac{(s_1-s_2)^2}{s_1+s_2},\qquad
 z=\frac{s_1-s_2}{s_1+s_2}.
\]
The function $J$ is nonnegative, convex, one-homogeneous and at most
$s_1+s_2$. Its derivatives are $J_1=2z-z^2$, $J_2=-2z-z^2$, hence
$|J_i|\leq3$, and
\begin{equation}\label{eq:JHess}
 D^2J=\frac{2w\otimes w}{s_1+s_2},\qquad w=(1-z,-1-z),\qquad
 \sum_i s_i|\partial_\alpha J_i|^2\leq
 2\,D^2J[\partial_\alpha s,\partial_\alpha s].
\end{equation}
The common conservative drift cancels by homogeneity. Both shifted equations
have source $\delta\Div_q(AqM)=-\delta M A:K$, whose contribution is at most
$C\delta\int|A|$.

Here are the required weak chain-rule details. First cut off configurations
at distance $r$ from the boundary and mollify in $(x,q)$ at radius $h<r/4$.
Rewriting the mollified equation with the original $b$ creates flux error
\[
 R_{i,h}(\xi)=\int[b(y)-b(\xi)]s_i(y)\eta_h(\xi-y)\dd y,\qquad
 \frac{|R_{i,h}|^2}{s_{i,h}}
 \leq\int|b(y)-b(\xi)|^2s_i(y)\eta_h(\xi-y)\dd y.
\]
Its integral on the fixed interior tends to zero: the $v$ term uses $L^2$
translation continuity and configuration mass $1+\delta$; for $Aq$ split
$(A(y_x)-A(x))y_q+A(x)(y_q-q)$ and use the same continuity and an
$O(h^2\int|A|^2)$ term. The field $\grad U$ is Lipschitz on that interior.
No spatial derivative of $A$ is taken. By \eqref{eq:JHess} the entropy error
is bounded by half the diffusion dissipation plus
$C/\min(\varepsilon,1)\sum_i|R_{i,h}|^2/s_{i,h}$.
The divergence error against the fixed cutoff vanishes as well.
Let $h\to0$ at fixed $r,\delta,\varepsilon$; retain the nonnegative
dissipation by lower semicontinuity. Bounded $J_i$, the $L^1$ diffusion
fluxes, and the bound on $\int|b|^2p_i$ justify the remaining passages by
drift truncation and time regularization.

Choose the configuration cutoff with derivative at most $C/d$ on
$d<2r$. The diffusion collar is bounded by
\[
 C\sum_i\left(\int_{d<2r}s_i/d^2\right)^{1/2}
             \left(\int_{d<2r}|\grad_q\sqrt{s_i}|^2\right)^{1/2}\to0.
\]
The confining collar is bounded by $C\int_{d<2r}\sum_i s_i/d^2\to0$;
the matrix-drift collar by
$C\sum_i(\int|A|^2s_i)^{1/2}(\int_{d<2r}s_i/d^2)^{1/2}\to0$.
All these integrals include time and space. Thus the common initial trace gives
\[
 \int J(p_1+\delta M,p_2+\delta M)(t)
       \leq C\delta\int_0^t\int_\T|A|.
\]
Dominated convergence as $\delta\to0$ yields $p_1=p_2$.
\end{proof}

To obtain the Fisher assumptions for $\E\rho$, approximate $v$ strongly in
$L^2_tH^1_x$ by smooth divergence-free $v_m$. For fixed $m$, $A_m$ is bounded;
\eqref{eq:fibreenergy}, volume preservation and Jensen put the averaged
$g_m=\E\rho_m/M$ in $L^\infty_tL^2_{x,q,M}$ with its configuration form bound.
Mollify its weak equation only in $x$. The $v_m$ and $A_m$ commutators tend
to zero in $L^2$, since their spatial derivatives are bounded and $g_m$ is
$L^2$. Testing by the spatially mollified $g_m$ gives the $\varepsilon$
spatial form bound. Removing this mollification and then the convex entropy
regularization gives
\begin{equation}\label{eq:averageFisher}
 \Ent(\E\rho_m(t))+2\int_0^t I_q(\E\rho_m)
 +4\varepsilon\int_0^t I_x(\E\rho_m)
 \leq\Ent(\psi_0)+C\int_0^t\int_\T|A_m|^2.
\end{equation}
Its constant is independent of $m$; no uniqueness assertion has yet been used.
The original bounded initial density is retained by the fibre $L^2$ energy
estimate when removing smooth initial approximations.

The flows $X_m$ converge in measure. Their pulled-back gradients converge
strongly in $L^2$, since
\[
 A_m(X_m)-A(X)=(A_m-A)(X_m)+A(X_m)-A(X);
\]
volume preservation controls the first term and bounded continuous
approximation of $A$ controls the second. The homogeneous entropy comparison
above, now with two matrix histories, absorbs their difference into the
configuration dissipation and gives on each fibre
\[
 \int_B J(Mr_m,Mr_\ell)(t)\dd q
 \leq C\int_0^t|A_m(X_m)-A_\ell(X_\ell)|^2\dd s.
\]
Its label integral and Cauchy--Schwarz imply strong $L^1$ convergence of the
fibre densities. Returning by the converging inverse flows identifies their
Eulerian limit. Lower semicontinuity in \eqref{eq:averageFisher} gives both
Fisher bounds for $\E\rho$. These also give its strong initial trace, either
by entropy or by comparison with the zero-matrix configuration semigroup.

The coupled solution has the same trace. In fact $\log g_0$ is bounded, weak
continuity and \eqref{eq:entropy} give
$\limsup_{t\downarrow0}\Ent(\psi(t))\leq\Ent(\psi_0)$, and hence
\[
 \operatorname{Ent}(\psi(t)\mid\psi_0)
   =\Ent(\psi(t))-\int\psi(t)\log g_0\longrightarrow0.
\]
Pinsker's inequality gives strong $L^1$ convergence. Lemma~\ref{lem:unique}
now proves $\psi_\varepsilon=\E\rho^\omega$. In particular the conditioning
has not changed the deterministic stress feedback in the fluid equation.

\subsection{Nonlinear transforms before expectation}\label{subsec:Ito}
If a fibre functional $B(t,b)$ has finite variation with source $S_B$, test
its Eulerian push-forward against $\phi(t,X(t,b))$. It\^o's product formula
has no quadratic covariation with $B$. Its sole correction is
$\varepsilon\int B\Delta_x\phi(t,X)\dd b$; its martingale coefficient is
$\sqrt{2\varepsilon}\int B\grad_x\phi(t,X)\dd b$.
For bounded functionals this coefficient is bounded. For $\Hone,\Htwo$
their deterministic global pathwise mass bounds make it square-integrable.
Fubini and expectation therefore give
\begin{equation}\label{eq:expectation}
 (\partial_t+v\cdot\grad_x)\E B
       =\varepsilon\Delta_x\E B+\E S_B.
\end{equation}
For $G$ first keep a fixed smooth configuration test and use the fibre
variational equation. Smooth fibre and spatial-flow approximations justify
the same identity; their fluxes and sources pass by the preceding estimates.
There is no spatial derivative of the conditional density in this argument.

For $\vartheta_\kappa(z)=z/(1+\kappa z)$ apply \eqref{eq:expectation} to
\begin{equation}\label{eq:quotient}
 B_{\delta,\kappa}=\frac{\vartheta_\kappa(\Hone)}{1+\delta\Htwo},\qquad
 W_{\delta,\kappa}=(1+\delta\Htwo)(1+\kappa\Hone)^2.
\end{equation}
The stress-work term is $A:\E[\tau(Ms)/W_{\delta,\kappa}]$, and the
configuration Fisher term has precisely the same conditional weight.
The denominator derivative is bounded in $L^1$ by $C\delta/\kappa$, using
\eqref{eq:logbounds}. Thus $\delta$ will be removed before $\kappa$.
The bound $B_{\delta,\kappa}\leq1/\kappa$ ensures that its spatial Laplacian
vanishes against fixed tests when $\varepsilon\to0$.
This identity concerns $\E B(\rho)$, not $B(\E\rho)$.
Applying the latter nonlinear functional directly to the spatially diffusive
density would introduce spatial Hessian terms with no required sign.

\section{A transported envelope for the limiting conditional moments}\label{sec:envelope}
Return to the coupled family in \eqref{eq:weakstart} and its conditional
representation. Along each random characteristic define
\begin{equation}\label{eq:Q}
 Q_n^\omega(t,X_n(t,b))=\mathcal H_{2,0}(b)+
                \int_0^t|F_n^\omega(s,X_n(s,b))|\dd s.
\end{equation}
It dominates $\mathcal H_{2,n}^\omega$, and its spatial mass is bounded by the same
deterministic constant on every path. For $\beta_L(z)=Lz/(L+z)$,
$q_{n,L}=\E\beta_L(Q_n^\omega)$ satisfies
\begin{equation}\label{eq:boundedQ}
 (\partial_t+u_n\cdot\grad_x)q_{n,L}
 =\varepsilon_n\Delta_xq_{n,L}
   +\E[\beta_L'(Q_n^\omega)|F_n^\omega|]\geq
      \varepsilon_n\Delta_xq_{n,L}.
\end{equation}
At fixed $L$ it is bounded, so strong velocity convergence passes its
transport and the small Laplacian vanishes. The limit $q_L$ is nondecreasing
along the regular Lagrangian flow of $u$. Its increasing supremum $Q$ is
finite almost everywhere by its common mass bound. Put
$R=\mathrm{w}\!\lim\E\mathcal H_{2,n}^\omega$. The $L^2$ bound on $\mathcal H_{2,n}^\omega$
and $\Htwo-\beta_L(\Htwo)\leq\Htwo^2/L$ show that $R\leq Q$.

On $[0,T_0]$ choose a later time $T_1>T_0$ at which $Q(T_1)$ has finite
mass, and transport $1+Q(T_1)$ backward by the limiting flow. This defines
\begin{equation}\label{eq:N}
 N\geq1,\qquad R\leq N,
 \qquad N(t,X(t,b))=N(0,b)<\infty\quad\text{a.e.}
\end{equation}
For this possibly unbounded $N$, transport means characteristic constancy
and transport of its bounded Borel functions. No raw product $uN$ is used.
The normalizer is introduced after the zero-diffusion limit; its later-time
definition is never an anticipating stochastic test.

Let $y=\mathrm{w}\!\lim\E G_n^\omega$. Weighted Cauchy--Schwarz yields
\[
 \frac{\int_BM|\grad_q\E G_n^\omega|^2}{\E\mathcal H_{2,n}^\omega}
 \leq\E\frac{\int_BM|\grad_qG_n^\omega|^2}{\mathcal H_{2,n}^\omega}.
\]
Convexity of $z\mapsto z\sqrt{\log z^2}$ on the shifted range and the $L^2$
triangle inequality also give
$\|\E G_n^\omega\sqrt{\log(\E G_n^\omega)^2}\|_{L^2_M}
 \leq C\E\mathcal H_{2,n}^\omega$.
Convex lower semicontinuity, \eqref{eq:logbounds}, and $R\leq N$ imply
\begin{equation}\label{eq:yenergy}
 \int_0^{T_0}\int_\T\frac{\Dy}{N}\leq C,
 \qquad \Dy=\int_BM|\grad_qy|^2,\qquad
 h:=y/N^2\in L^\infty_{t,x}L^2_M\cap L^2_{t,x}H^1_M.
\end{equation}
These are the energy bounds needed later for deterministic configuration
renormalization. Only bounded envelope truncations have entered a spatial
transport limit.

\section{Conditional cutoffs and the stress-work defect}\label{sec:work}
The stress-work identity must use the same conditional cutoffs as the entropy
equation. For $c=(\delta,\kappa)$, define
\begin{equation}\label{eq:cutstress}
 T_{n,c}=\E\frac{\tau(Ms_n^\omega)}{W_{n,c}^\omega},\qquad
 z_{n,c}=A_n\E\frac{s_n^\omega}{W_{n,c}^\omega},\qquad
 W_{n,c}^\omega=(1+\delta\mathcal H_{2,n}^\omega)(1+\kappa\mathcal H_{1,n}^\omega)^2.
\end{equation}
The raw shifted stress is $\tau_n^a=\tau(\psi_n)+a\Id$.
All cutoff removals below first take $n\to\infty$ at fixed $c$, then
$\delta\downarrow0$, and then $\kappa\downarrow0$.

\subsection{Integrable stress squares and a common drag direction}\label{subsec:cutbounds}
For a conditional profile let
$S=\int_BMs/d$, $J_q=\int_BMs\log s/d$ and
$D_G=\int_BM|\grad_qG|^2$. Applying \eqref{eq:Hardy} to $G$ gives
\[
 J_q^2\leq C\Hone(D_G+\Hone).
\]
Jensen's inequality for the finite measure $M\dd q/d$ gives
$S^2\log^2(2+S)\leq C(1+J_q^2)$. Since $\Hone\leq C\Htwo$,
\begin{align}
 \left|\frac{\tau(Ms)}{W_c}\right|^2
 \log^2\left(2+\left|\frac{\tau(Ms)}{W_c}\right|\right)
 &\leq C_\delta(1+D_G/\Htwo),\label{eq:stressOrlicz}\\
 \int_B\frac M d\,\Phi(|A|s/W_c)
 &\leq C_\delta(1+|A|^2+D_G/\Htwo),\qquad
 \Phi(z)=z\log(1+z).\label{eq:dragOrlicz}
\end{align}
For the second inequality one also uses
$bc\log(2+b)\leq C[b^2+c^2\log^2(2+c)+1]$ for $b,c\geq0$.
The right sides have uniformly bounded integrals by \eqref{eq:logbounds}.
The superlinear factors on the left therefore give uniform integrability.
Jensen transfers it to $|T_{n,c}|^2$ in $(t,x)$ and to $z_{n,c}$ in
$M\dd q\dd x\dd t/d$.

The conditional raw stress is bounded in joint $L^2$. The cutoff factors
approach one in measure uniformly in $n$ by the $L^1$ moment bounds.
Consequently
\begin{equation}\label{eq:Lpcut}
 \sup_n\|T_{n,c}-\tau_n^a\|_{L^p_{t,x}}\longrightarrow0
       \qquad(1<p<2).
\end{equation}
The radial matrix $K$ is positive semidefinite. The $T_{n,c}$ increase in the
Loewner order when their scalar cutoffs are removed, and their squared
Frobenius norms increase as well, since
$|S+D|^2-|S|^2=2\tr(SD)+|D|^2\geq0$ for positive semidefinite $S,D$.
At fixed $(t,x,q)$, all $z_{n,c}$ have the common matrix direction $A_n$ and
nonnegative ordered scalar magnitudes. Thus for two ordered cutoffs
\begin{equation}\label{eq:polar}
 |z_{n,c}-z_{n,c'}|=|z_{n,c}|-|z_{n,c'}|
 \quad\text{when }W_{n,c}\leq W_{n,c'}.
\end{equation}
Expectation has preserved both monotonicities needed here.

For clarity, this identifies a weighted limit of the raw drag, not merely a
cutoff-dependent product. The raw $A_ns_n$ are uniformly integrable in the
unweighted $M$ measure: the superlinear function
$\Phi_{1/2}(z)=z\sqrt{\log(1+z)}$ obeys
\[
 \Phi_{1/2}(|A|s)\leq Cs(1+|A|^2+\log s).
\]
Normalization makes its right side integrable after averaging. Also, by
\eqref{eq:Hardy},
\[
 \int M|A_n|\E s_n/d\leq C\int(1+|A_n|^2+\E I_q(\rho_n)).
\]
At fixed cutoffs extract weighted weak
$L^1$ limits $z_c$ and $g_c$ of $z_{n,c}$ and $|z_{n,c}|$.
Equation~\eqref{eq:polar} implies
$|z_c-z_{c'}|\leq g_c-g_{c'}$. The increasing $g_c$ have bounded weighted
mass, so $z_c$ are Cauchy in weighted $L^1$. Their limit coincides with the
raw weak product by the unweighted uniform integrability. Denote
\begin{equation}\label{eq:beta}
 \beta=\mathrm{w}\!\lim(A_n\psi_n)-A\psi.
\end{equation}
The work covariance is therefore exactly $\int_BK:\beta$, with no second
product selected by removing the cutoffs.

\subsection{Periodic linear estimates and the compact convective part}\label{subsec:periodicStokes}
Split $u_n=v_n+w_n$. Here $v_n$ has zero initial data and solves the periodic
solenoidal Stokes equation with force $\Div_x\tau(\psi_n)$; $w_n$ has initial
data $u_0$ and force $-u_n\cdot\grad_xu_n$. Let $v_{n,c}$ be the analogous
zero-initial Stokes solution forced by $\Div_x T_{n,c}$.
We need the linear $L^p$ gradient estimate for stress forcing and the
second-derivative estimate for vector forcing, $1<p<2$, used in
\cite[Section~5.4]{Masmoudi}. We record why their spatially periodic versions
hold.

The causal kernels in both estimates have the form $D_x^2\mathbb P G_\nu$,
where $G_\nu$ is the heat kernel and $\mathbb P$ is the solenoidal projector.
They are sums of second Gaussian derivatives and fourth derivatives of the
Newton-potential convolution with a Gaussian. For $0<t<T$ and sufficiently
large $|x|$ their absolute values are at most $C_{\nu,T}|x|^{-5}$.
To see this, differentiate the Newton kernel four times and split its
convolution at $|y|=|x|/2$; the remaining Gaussian tail has faster decay.
Thus the space-time tails of these kernels are integrable and tend to zero
outside large spatial balls.

Cut a smooth periodic forcing, extended to $\R^3$, off to a cube $Q_L$.
On its interior at distance $\sqrt L$ from the boundary, the whole-space and
periodic convolutions differ by at most the forcing supremum times the kernel
tail. Apply the whole-space $L^p$ estimate, divide its $p$th-power bound by the
number of periodic cells, and let $L\to\infty$. The interior fraction tends
to one and the error tends to zero. Density extends the resulting periodic
bound to all $L^p$ forcings. This transfers a linear estimate only; it does not
transfer a nonlinear existence theorem between domains.

The convective force belongs to $L^{5/4}_{t,x}$ by the fluid energy bounds.
The transferred estimate gives $w_n\in L^{5/4}_tW^{2,5/4}_x$.
Its mean-zero pressure is $-(-\Delta)^{-1}\Div(-u_n\cdot\grad u_n)$;
the periodic $|x|^{-2}$ kernel is integrable, giving an $L^{5/4}$ bound.
Hence $\partial_tw_n$ is bounded in $L^{5/4}_tW^{-1,5/4}_x$.
Spatial smoothing and finite-dimensional time compactness give compactness
of $\grad w_n$ in $L^{5/4}_{t,x}$; interpolation with its $L^2$ bound gives
\begin{equation}\label{eq:wcompact}
 \grad w_n\longrightarrow\grad w\quad\text{in }L^p_{t,x},\qquad p<2.
\end{equation}
The fixed initial-data Stokes part is common to all $n$. From
\eqref{eq:Lpcut} and the stress-to-gradient estimate,
$\grad v_{n,c}-\grad v_n\to0$ in $L^p$, uniformly in $n$, as $c$ is removed.

\subsection{Local Stokes work and retained oscillation}\label{subsec:localwork}
At fixed $c$, write $T_c=\mathrm{w}\!\lim T_{n,c}$ and
$G_c=\mathrm{w}\!\lim|T_{n,c}|^2$. By \eqref{eq:stressOrlicz}, $G_c$ is
$L^1$. The $G_c$ increase and have bounded masses, so $G_c\uparrow G$ in
$L^1$ for a finite $G$. The needed tail comparison concerns these weak
limits. It does not assert uniform integrability of the raw stress squares.
Indeed use the monotone caps $\min(|T_{n,c}|^2,L^2)$. Their weak limits
increase both with $L$ and with cutoff removal; the two suprema commute and
equal $G$. Given $\zeta>0$, choose $c_0,L_0$ so that the gap from $G$ to the
doubly capped weak limit has $L^1$ norm less than $\zeta$. For every smaller
$c$, the weak limiting energy in $|T_{n,c}|>2L_0$ is at most
$4\zeta/3$, since
\[
 b^2\mathbf1_{b>2L}\leq\tfrac43\{b^2-\min(b^2,L^2)\}.
\]

Pair a bounded capped stress with the $L^p$-small difference of the cut and
uncut Stokes gradients, then remove the cap with this tail comparison and
the uniform $L^2$ gradient bounds. It follows that the weak work limits
$T_{n,c}:\grad v_{n,c}$ and $T_{n,c}:\grad v_n$ have the same $L^1$ limit
as $c$ is removed. Equation~\eqref{eq:wcompact} and fixed-$c$ uniform
integrability also give
$\mathrm{w}\!\lim(T_{n,c}:\grad w_n)=T_c:\grad w$.

On the other hand $0\leq T_c\leq\tau(\psi)+a\Id$ in the Loewner order,
and \eqref{eq:Lpcut} gives $T_c\to\tau(\psi)+a\Id$ in $L^p$, $p<2$.
Domination by this $L^2$ stress gives strong $L^2$ convergence, hence strong
$L^2$ convergence of the limiting cutoff Stokes gradients.

For fixed $c$, test the Stokes equation by its velocity times a smooth
space-time cutoff. Its energy-class velocity is strongly continuous in
$L^2$; pressure and force fluxes pass by strong velocity and weak
pressure/stress convergence. Comparing the local energy identities gives
\begin{equation}\label{eq:localenergy}
 \mathrm{w}\!\lim(T_{n,c}:\grad v_{n,c})
       =T_c:\grad v_c-\nu\mu_c,
 \qquad
 \mu_c=\mathrm{w}\!\lim|\grad v_{n,c}|^2-|\grad v_c|^2\geq0.
\end{equation}
Remove the cutoffs by the preceding capped-work comparison and the
common-direction argument of \eqref{eq:polar}. The total work limit is
$A:(\tau(\psi)+a\Id)+\int_BK:\beta$ by \eqref{eq:beta}.
Subtract the compact $w$ work from \eqref{eq:localenergy}. We obtain
\begin{equation}\label{eq:mu}
 \nu\mu=-\int_BK:\beta,\qquad \mu\geq0,\qquad \mu\in L^1_{t,x}.
\end{equation}
The last assertion follows because the resulting work terms are $L^1$
functions. The isotropic shift contracts to zero with the trace-free $A$.
This is the local-energy comparison of \cite[Proposition~5.4]{Masmoudi},
with the conditional PSD, cap, and common-direction properties verified above.

We will also use its oscillation consequence. If $\Pi$ is the velocity-gradient
marginal of the joint law constructed below, then
\begin{equation}\label{eq:PiVar}
 \pair{|\Pi-A|^2}\leq\mu.
\end{equation}
Here is a bounded-moment way to check this consequence. For a centered gradient
matrix use $P_R(H)=\min(|H|^2,R)$, whose Lipschitz constant is at most
$2\sqrt R$. The $L^p$-small cut-gradient difference makes their $P_R$ moments
coincide in the cutoff limit. The cut-gradient square defect bounds its
centered $P_R$ moment. Remove the cutoff first, then let $R\to\infty$.
The compact $\grad w_n$ adds no oscillation. This proves
\eqref{eq:PiVar}; it does not identify $\mu$ with the entire viscous
dissipation defect of $u_n$.

\section{Joint conditional fluctuations and configuration coercivity}\label{sec:law}
We now compare the weak limit of the conditional entropy with the square of
its weakly limiting square root. Define
\begin{equation}\label{eq:defect}
 f=\mathrm{w}\!\lim\E f_n^\omega,\quad
 y=\mathrm{w}\!\lim\E G_n^\omega,\quad
 \theta=\mathrm{w}\!\lim\E\mathcal H_{1,n}^\omega,\quad
 \eta=\theta-\int_BMy^2\geq0.
\end{equation}
In particular $\eta\leq CN$ by \eqref{eq:N} and $\Hone\leq C\Htwo$.
To retain correlations with $A_n$, construct at each $(t,x,q)$ the joint law
of $(A_n,f_n^\omega,\grad_qf_n^\omega)$ before averaging the conditional
variable. One can do this without a Young-measure theorem on Brownian path
space: extract weak-star limits of
$\E\Phi(A_n,f_n^\omega,\grad_qf_n^\omega)$ for a countable dense set of
compactly supported continuous state tests. Positivity gives a positive
functional on the finite-dimensional state space. Riesz representation gives
a measure; the moment bounds give tightness and total mass one. Denote its
state by $(\Pi,\lambda,\Lambda)$ and its expectation by brackets.
The $\Pi$ marginal is independent of $q$, is trace-free, has mean $A$, and
satisfies \eqref{eq:PiVar}.

The higher moment of $f_n$ and the $L^2$ gradient bound make $\lambda\Lambda$
uniformly integrable. They likewise absorb the logarithmic multiplier in
the gradient of $G$. Consequently, with $b=\sqrt{\log\lambda^2}$,
\begin{align}
 f&=\pair\lambda,&\grad_qf&=\pair\Lambda,&
 y&=\pair{\lambda b},&\grad_qy&=\pair{(b+b^{-1})\Lambda},\label{eq:lawmoments}\\
 \overline s&:=\pair{\lambda^2}=\psi/M+a.
 \nonumber
\end{align}
The square-root variance $V_f=\pair{\lambda^2}-f^2$ obeys
\begin{equation}\label{eq:variancegrad}
 \grad_qV_f=2\pair{(\lambda-f)(\Lambda-\grad_qf)},\qquad
 |\grad_q\sqrt{V_f}|^2\leq\pair{|\Lambda-\grad_qf|^2}.
\end{equation}
Since $\lambda\mapsto\lambda\sqrt{\log\lambda^2}$ has derivative bounded
below on $\lambda\geq\sqrt a$, the two-copy variance formula gives
$\int_BMV_f\leq C\eta$.

\subsection{Retaining positive Fisher moments}\label{subsec:retention}
At each fixed conditional entropy cutoff, the positive measures
$\E[|\grad_qf_n^\omega|^2/W_{n,c}^\omega]$ have weak limits. They increase
when the cutoffs are removed and have bounded total masses. Let $m$ be the
increasing limit of their absolutely continuous densities. For a bounded
state moment $\min(|\Lambda|^2,L^2)$, the cutoff loss is at most
\[
 L^2\E(1-1/W_{n,c})
 \leq L^2\E\min(\delta\mathcal H_{2,n}+2\kappa\mathcal H_{1,n},1),
\]
whose total integral tends to zero as $c$ is removed. Weak passage of this
bounded moment, then removal of $c$ and $L\to\infty$, gives
\begin{equation}\label{eq:mretain}
 m\geq\pair{|\Lambda|^2}\quad\text{a.e.},\qquad
 \varpi:=\int_BM(m-|\grad_qf|^2)\geq0.
\end{equation}
This is a density inequality: for an increasing sequence of finite positive
measures, the absolutely continuous part of the limit is the increasing limit
of their absolutely continuous parts, since all their singular parts are
supported on a countable union of null sets.

The positive transformed source is
$-\Theta''(s_n)|\grad_qs_n|^2$. Its coefficient in the $f_n$ variables is
$(b+b^{-3})/\lambda$, bounded and continuous on $\lambda\geq\sqrt a$.
The same bounded-state argument implies that its absolutely continuous limit
$V$ satisfies
\begin{equation}\label{eq:Vretain}
 V\geq\pair{(b+b^{-3})|\Lambda|^2/\lambda}.
\end{equation}
The singular positive part remains favorable in a supersolution inequality;
we will discard it before using a rough spatial normalizer.
From \eqref{eq:variancegrad}, \eqref{eq:mretain}, and Hardy,
\begin{equation}\label{eq:varianceHardy}
 \int_B\frac{MV_f}{d}\leq C\sqrt\eta\sqrt{\varpi+\eta}.
\end{equation}

\subsection{An all-state two-copy inequality}\label{subsec:matrix}
The following algebra ensures that averaging has not destroyed configuration
coercivity. Take two independent states $(\lambda,\Lambda)$,
$(\lambda',\Lambda')$ of the same law, and put
$b=\sqrt{\log\lambda^2}$, $c=\sqrt{\log(\lambda')^2}$, $r=\lambda'/\lambda$.
Set
\[
 B_1=2+rc(b+b^{-3}),\quad
 B_2=2+r^{-1}b(c+c^{-3}),\quad
 B_3=(b+b^{-1})(c+c^{-1}).
\]
The matrix with diagonal $B_1-1,B_2-1$ and off-diagonal $-(B_3-1)$ is
positive semidefinite. Its diagonal entries are positive, and direct expansion
of its determinant gives
\begin{align*}
 &bc[r+r^{-1}+2-2b/c-2c/b]+2[b/c+c/b-2]\\
 &\hspace{12mm}+(bc)^{-1}
       [rc^2/b^2+r^{-1}b^2/c^2+2-2b/c-2c/b].
\end{align*}
By symmetry take $r\geq1$, $h=\log r$, and
$z=\tfrac12\log(c^2/b^2)$. Since $b^2\geq4$ and $c^2-b^2=2h$,
$0\leq z\leq h/2$. The first bracket is
$2\cosh h+2-4\cosh z\geq0$, the middle bracket is nonnegative, and the
last is $2\cosh(h+2z)+2-4\cosh z\geq0$.
It follows, for every pair of admissible states, that
\begin{equation}\label{eq:pairmatrix}
 B_1|\Lambda|^2+B_2|\Lambda'|^2-2B_3\Lambda\cdot\Lambda'
       \geq|\Lambda-\Lambda'|^2.
\end{equation}
Take the expectation over the two independent copies. The mixed terms factor,
using $y=\pair{\lambda b}$ and
$\grad_qy=\pair{(b+b^{-1})\Lambda}$, to give the pointwise inequality
\begin{equation}\label{eq:pairaverage}
\begin{aligned}
 4\pair{|\Lambda|^2}
 +2y\pair{(b+b^{-3})|\Lambda|^2/\lambda}
 -2|\grad_qy|^2
 \geq 2\bigl(\pair{|\Lambda|^2}-|\grad_qf|^2\bigr).
\end{aligned}
\end{equation}
Replace the two positive moments on the left by $m$ and $V$ using
\eqref{eq:mretain}--\eqref{eq:Vretain}. After subtracting
$2(m-|\grad_qf|^2)$ from the new left side, the retained Fisher excess leaves
at least $2(m-\pair{|\Lambda|^2})\geq0$. Integration against $M$ gives
\begin{equation}\label{eq:Fishercoercive}
 4\int_BMm+2\int_BMyV-2\Dy\geq2\varpi.
\end{equation}
The joint-law gradient variance contributes at least twice its mass.
The excess $m-\pair{|\Lambda|^2}$ contributes four times its mass, so
retained concentration has the favorable sign without being assumed zero.

\subsection{Covariances and configuration boundary integrability}\label{subsec:covariances}
Use $s=\lambda^2$ within law brackets and put
\begin{equation}\label{eq:Gammas}
 L=\pair{s\Theta'(s)},\quad P=\pair{\Theta'(s)},\quad
 \Gamma=\pair{\Pi G}-Ay,\quad
 \Gamma'=\pair{\Pi s\Theta'(s)}-AL.
\end{equation}
Cauchy--Schwarz and \eqref{eq:PiVar} give
\begin{equation}\label{eq:basiccov}
 |\Gamma|^2\leq\mu\Var(G),\qquad
 |\Gamma'-\Gamma|^2\leq\tfrac14\mu\Var(G).
\end{equation}
For the second bound note that $s\Theta'(s)-G$, as a function of $G$, has
derivative of absolute value at most $1/2$. Since
$\int_BM\Var(G)=\eta\leq CN$, these imply, before any renormalization,
\begin{equation}\label{eq:normalizedcov}
 \int_BM|\Gamma/N^2|^2+
 \int_BM|\Gamma'/N^2|^2\leq C\mu/N^3.
\end{equation}
Thus the normalized covariance fluxes belong to $L^2(M)$ and, by
Cauchy--Schwarz with $M/d^2$, to $L^1(M/d)$. The mean-gradient flux
$Aqh$ is $L^2(M)$ by the uniform $q$ norm in \eqref{eq:yenergy}.

For signed scalar sources a direct profile estimate uses exactly $k>2$.
Choose $2/3<\alpha<k/(k+1)$. Then $G^{4/3}\leq C_{a,\alpha}s^\alpha$ and
\begin{equation}\label{eq:profile}
 \int_B\frac M dG^{4/3}
 \leq C\left(\int_BMs\right)^\alpha
          \left(\int_BM d^{-1/(1-\alpha)}\right)^{1-\alpha}
 \leq C_{k,a}.
\end{equation}
Hence $A_nG_n$, $A_ns_n\Theta'(s_n)$, and the lower shifted composites
are bounded in the mixed space
\[
 L^2_{t,x}\bigl(L^{4/3}(B,M\dd q/d)\bigr),
\]
including the conditional variable. Their signed limits have no concentration
in that measure.
For the conditional diffusion collar, Cauchy--Schwarz gives
\[
 \int_{t,x,d<r}\frac M d|\grad_qG_n|
 \leq\left(\int\frac{M|\grad_qG_n|^2}{\mathcal H_{2,n}}\right)^{1/2}
       \left(\int_{t,x}\mathcal H_{2,n}\int_{d<r}\frac M{d^2}\right)^{1/2}\to0.
\]
The transformed-drag collar is at most
$C\|A_n\|_{L^1_{t,x}}(\int_{d<r}M/d)^{1/4}\to0$ by \eqref{eq:profile}.
Both vanish uniformly. The normalized limiting fluxes have the same boundary
property by \eqref{eq:yenergy} and \eqref{eq:normalizedcov}. Consequently
the transformed variational inequalities hold for tests meeting the closed
configuration ball, not only for interior tests.

\section{The normalized entropy-defect inequality}\label{sec:defect}
We next justify the limiting transport and configuration operations before
estimating their four resulting terms. This order matters because the raw
entropy moment $\theta$ need not have spatial integrability above exponent two.

\subsection{The lower transformed inequality}\label{subsec:lower}
The averaged $G$ equation has just the linear spatial diffusion term in
\eqref{eq:expectation}. At the limit its positive configuration source is a
measure. Discard its singular positive part, retaining the density $V$ of
\eqref{eq:Vretain}. Write $D_t=\partial_t+u\cdot\grad_x$ and
$\mathcal L_Mh=M^{-1}\Div_q(M\grad_qh)$. The resulting inequality is
\begin{equation}\label{eq:ylower}
 D_ty\geq\mathcal L_My
 -\frac1M\Div_q\bigl(M\pair{\Pi G}q\bigr)
 +K:\pair{\Pi[(s-a)\Theta'(s)-G]}+V.
\end{equation}
Here $s=\lambda^2$, $G=\Theta(s)=\lambda\sqrt{\log\lambda^2}$ inside the
brackets, $\pair{\Pi G}=Ay+\Gamma$, and
$\pair{\Pi s\Theta'(s)}=AL+\Gamma'$.
For a single fibre with its trace-free matrix history $\mathsf B(t)$, the
chain rule giving these signs is
\[
 \partial_tG=\mathcal L_MG
 -\frac1M\Div_q(M\mathsf BqG)
 +(\mathsf B:K)[(s-a)\Theta'(s)-G]
 -\Theta''(s)|\grad_qs|^2.
\]
Use \eqref{eq:expectation} before taking the joint-law limit; the last source
is positive and has retained density $V$. The configuration boundary inequality
and lower initial trace $y(0)\geq\Theta(g_0+a)$ are inherited in the
variational sense. Thus \eqref{eq:ylower} does not impose a pointwise boundary
value or assert equality of a residual measure.
Its scalar pairings with fixed nonnegative smooth configuration tests have
$L^r_{t,x}$ regularity for some $r>2$. To interpret these along the limiting
spatial flow, mollify in $x$. The Sobolev transport commutator tends to zero
in $L^1$, since $\grad u\in L^2$ and the scalar pairing has exponent above
two. Apply an increasing bounded scalar truncation, use the bounded mild
transport inequality of \cite[Proposition~A.5]{Masmoudi}, then remove the
truncation by dominated convergence for the signed $L^1$ source and monotone
convergence for the positive residual. A countable dense family of
nonnegative configuration tests and the flux integrability give a common
full-measure set of labels on which the whole configuration inequality holds.

Along each such label $N$ is a fixed finite scalar. Division gives
\begin{equation}\label{eq:hnormalized}
\begin{aligned}
 D_th\geq{}&\mathcal L_Mh
 -\frac1M\Div_q\bigl(M(Ah+\Gamma/N^2)q\bigr)\\
 &+\frac1{N^2}K:\pair{\Pi[(s-a)\Theta'(s)-G]}+\frac V{N^2}.
\end{aligned}
\end{equation}
This is the labelwise variational inequality for $h=y/N^2$.
By \eqref{eq:yenergy}, \eqref{eq:normalizedcov}, and
\eqref{eq:profile}, $h=y/N^2$ belongs to
$L^\infty_tL^2_M\cap L^2_tH^1_M$, its force flux is both $L^2(M)$ and
$L^1(M/d)$, and its signed scalar source and $V/N^2$ are $L^1(M)$.
The boundary inequality and initial trace are inherited from the conditional
variational identities. These are the hypotheses of
\cite[Proposition~B.6]{Masmoudi}, now for a deterministic configuration
supersolution on a label. In its ordinary-divergence notation, rewrite the
weighted force divergence as ordinary divergence plus the $\grad U$ term;
the $L^1(M/d)$ flux bound controls that additional signed source. We use the
convex supersolution statement of B.6 here, with its lower initial and boundary
inequalities, and do not invoke an equality-class continuity assertion.

Apply smooth convex approximations to the square with bounded nonnegative
derivative. Their flux cross products are integrable by the normalized
$L^2$ bounds. For the signed sources note that $L\leq Cy$ and $P$ is bounded.
Quadratic Hardy and \eqref{eq:basiccov} bound their products with $y/N^4$ by
\begin{equation}\label{eq:squareproducts}
 C\left[\frac{1+|A|^2}{N^3}+\frac{\Dy}{N^4}+\frac\mu{N^3}\right].
\end{equation}
These are integrable by \eqref{eq:yenergy}. Remove the derivative cap by
dominated convergence for signed terms and monotone convergence for $yV$.
No fourth configuration moment of $h$ has been used. Because the singular
positive measure was discarded first, no undefined representative of $N^{-2}$
has been multiplied into such a measure.

\subsection{The upper entropy inequality}\label{subsec:upper}
Keep the bounded quotients \eqref{eq:quotient} until after normalization.
At each fixed $\delta,\kappa$, their limiting equation gives an upper
transport inequality for a bounded scalar, with $L^1$ signed work and a
nonpositive Fisher residual. Discard only the singular negative part of that
residual, apply the bounded mild transport inequality, and multiply along the
characteristic by its fixed scalar $N(0,b)^{-4}$. Only then let
$\delta\downarrow0$ and $\kappa\downarrow0$.

The limiting quotient fields converge strongly in $L^1$ to $\theta$ in this
cutoff limit: their difference from the prelimit $\mathcal H_{1,n}$ is uniformly
small by joint uniform integrability. The signed work sources converge strongly
in weighted $L^1$ by \eqref{eq:polar} and the cutoff identification following
it. The shifted signed terms have a uniform exponent above one by
\eqref{eq:profile}. The denominator error is $O(\delta/\kappa)$.
The retained absolutely continuous Fisher densities increase to $m$, so
monotone convergence remains valid after multiplication by $N^{-4}$.
The initial quotients converge to $\mathcal H_{1,0}/N(0,b)^4$.
This proves the required upper inequality for $\theta/N^4$ without a spatial
commutator for raw $\theta$ or a raw product $u\theta$.
Its complete source expression, along the characteristic, is
\begin{equation}\label{eq:thetaupper}
\begin{aligned}
 D_t(\theta/N^4)\leq\frac1{N^4}\biggl[
 &-4\int_BMm+A:\tau(M\overline s)-\nu\mu\\
 &-a\int_BMK:\pair{\Pi(\log s+1)}\biggr].
\end{aligned}
\end{equation}
The raw conditional entropy source is
$-4\int M|\grad_qf_n|^2+A_n:\tau(Ms_n)
 -a\int M(A_n:K)\log s_n$.
The work term becomes the displayed mean work minus $\nu\mu$ by
\eqref{eq:mu}. Adding $1$ inside the last logarithmic bracket changes its
integral by zero, since $\int_BMK=\Id$ and $\tr A=0$.
Equation~\eqref{eq:thetaupper} is obtained only after the bounded-quotient
normalization and ordered cutoff limits above.

\subsection{The four terms and their signs}\label{subsec:four}
Subtract the lower squared-transform inequality from the upper entropy
inequality. Put
\begin{equation}\label{eq:e}
 e=\frac\eta{N^4}=\frac\theta{N^4}-\int_BMh^2,\qquad
 C_1=\pair{\lambda^2}-y\pair{\lambda/\sqrt{\log\lambda^2}}.
\end{equation}
The square passage from \eqref{eq:hnormalized} gives the complementary balance
\begin{equation}\label{eq:hsquare}
\begin{aligned}
 D_t\int_BMh^2\geq\frac1{N^4}\biggl[
 &-2\Dy+2\int_BM(\pair{\Pi G}q)\cdot\grad_qy\\
 &+2\int_BMyK:\pair{\Pi[(s-a)\Theta'(s)-G]}
   +2\int_BMyV\biggr].
\end{aligned}
\end{equation}
For the mean flux, weighted integration by parts and $\tr A=0$ give
$2\int_BM(Ayq)\cdot\grad_qy=\int_BM(A:K)y^2$.
Also $2L=y+\pair{\lambda/\sqrt{\log\lambda^2}}$, so
$\overline s-2yL+y^2=C_1$.
Insert $\pair{\Pi G}=Ay+\Gamma$ and
$\pair{\Pi s\Theta'(s)}=AL+\Gamma'$ in \eqref{eq:hsquare}, and subtract it
from \eqref{eq:thetaupper}. The mean, Fisher, shifted, and covariance terms
are, respectively,
\begin{equation}\label{eq:fouridentity}
 (\partial_t+u\cdot\grad_x)e\leq
       -\frac{\mathcal B_1+\mathcal B_2+\mathcal B_3+\mathcal B_4}{N^4},
\end{equation}
where, with $s=\lambda^2$ inside brackets,
\begin{align}
 \mathcal B_1&=\nu\mu-A:\tau(MC_1),\label{eq:B1}\\
 \mathcal B_2&=4\int_BMm+2\int_BMyV-2\Dy,\label{eq:B2}\\
 \mathcal B_3&=a\int_BMK:
       \pair{\Pi(\log s+1)-2y\Pi\Theta'(s)},\label{eq:B3}\\
 \mathcal B_4&=2\int_BM\{(\Gamma q)\cdot\grad_qy
                               +(\Gamma'-\Gamma):K\,y\}.\label{eq:B4}
\end{align}
In \eqref{eq:B1}, $\tau(MC_1)=\int_BKMC_1$.
The unnormalized expressions in \eqref{eq:B1}--\eqref{eq:B4} specify the
integrands in the normalized weak identity \eqref{eq:fouridentity}; no raw
transport equation for an unbounded $\eta$ is asserted.
The stress covariance is \eqref{eq:mu}; the remaining terms are precisely the
mean stress mismatch, configuration chain-rule dissipation, shift, and
gradient-density covariance when the two equations are subtracted.

Both maps $\lambda\mapsto\lambda\sqrt{\log\lambda^2}$ and
$\lambda\mapsto\lambda/\sqrt{\log\lambda^2}$ are increasing on the shifted
range. Their covariance gives $0\leq C_1\leq V_f$ by Cauchy--Schwarz.
Thus \eqref{eq:varianceHardy} and Young's inequality imply
\begin{equation}\label{eq:B12bounds}
 \mathcal B_1\geq\nu\mu-\tfrac18\varpi-C(1+|A|^2)\eta,
 \qquad \mathcal B_2\geq2\varpi.
\end{equation}
The second assertion is \eqref{eq:Fishercoercive}.

For the shift, use $\log s+1=2\Theta'(s)G$ and split $\Pi=A+(\Pi-A)$.
The fluctuating part is at most $C\sqrt{\mu\eta}$ because $\Theta'$ is
bounded on $s\geq a$ and $\int M/d^2<\infty$. The mean-gradient part is
bounded by $C|A|\int MV_f/d$. The scalar estimate behind this bound is
\begin{equation}\label{eq:shiftpair}
 |[\Theta'(\lambda^2)-\Theta'((\lambda')^2)]
   [\lambda\sqrt{\log\lambda^2}-\lambda'\sqrt{\log(\lambda')^2}]|
       \leq C_a|\lambda-\lambda'|^2.
\end{equation}
For $\lambda'/\lambda\leq2$ it follows by the displayed scalar derivatives.
For a larger ratio use
$\Theta'(\lambda^2)\leq C\sqrt{\log\lambda^2}/\lambda$ and
$|\lambda'-\lambda|\geq\lambda'/2$; the remaining factor is bounded by
$C\log(\lambda^2)/\lambda^2$. Symmetry covers the reversed ratio.
Therefore
\begin{equation}\label{eq:B3bound}
 |\mathcal B_3|\leq\tfrac\nu8\mu+\tfrac18\varpi
                                     +C(1+|A|^2)\eta.
\end{equation}
Finally \eqref{eq:basiccov}, Cauchy--Schwarz in $q$, and Hardy for $y$ give
\begin{equation}\label{eq:B4bound}
 |\mathcal B_4|\leq\tfrac\nu8\mu+C\eta(\Dy+N).
\end{equation}
All constants may depend on $\nu,k,a,T$ and the datum, with no smallness
condition on these data and no positive lower bound on $\nu$ beyond $\nu>0$.

\subsection{Characteristic closure}\label{subsec:Gronwall}
Combining \eqref{eq:fouridentity}--\eqref{eq:B4bound} and dropping the
remaining favorable terms gives
\begin{equation}\label{eq:Gronwall}
 (\partial_t+u\cdot\grad_x)e
 \leq C(1+|A|^2+N+\Dy)e,
 \qquad 0\leq e\leq C/N^3.
\end{equation}
The right side is globally $L^1$: use $N\geq1$, $A\in L^2$, and
$\Dy/N\in L^1$. Along almost every limiting characteristic, $N$ is a
finite constant, $\int|A|^2<\infty$, and $\int\Dy/N<\infty$ by volume
preservation. The coefficient of $e$ is therefore integrable in time on that
characteristic. The upper entropy and lower transform initial traces give
$e(0)\leq0$, hence $e(0)=0$. Applying the bounded mild transport inequality
and the ordinary characteristic Gronwall inequality yields
\begin{equation}\label{eq:eta0}
 \eta=0\quad\text{on every finite interval.}
\end{equation}
This closes the conditional density fluctuation without any strong-gradient
premise.

\section{Strong density convergence and the original weak system}\label{sec:final}
The definition of $\eta$, its vanishing, and uniform integrability of the
conditional $G_n^2$ give
\begin{equation}\label{eq:strongG}
 \int_0^T\int_\T\int_BM\E|G_n^\omega-y|^2\longrightarrow0.
\end{equation}
On $s,t\geq a$ the inverse transform satisfies
$|s-t|\leq C(1+\Theta(s)+\Theta(t))|\Theta(s)-\Theta(t)|$.
Cauchy--Schwarz thus turns \eqref{eq:strongG} into strong $L^1$ convergence
of the conditional densities to $M(\Theta^{-1}(y)-a)$.
Their averages converge strongly to the same density, which must be the weak
limit $\psi$. In particular every bounded kinetic test moment converges
strongly in $L^2_{t,x}$: its difference is bounded by a constant times the
conditional total variation, itself at most two, and tends to zero in $L^1$.
Pairing with $A_n\weak A$ in $L^2$ proves convergence of the exact
configuration drag.

Strong density convergence identifies stress truncations strongly; the uniform
collars \eqref{eq:tails} then give \eqref{eq:tauconv}.
Spatial transport passes as in Section~\ref{sec:prelim}.
Moreover
\[
 \int M|\sqrt{\psi_n/M}-\sqrt{\psi/M}|^2\leq\int|\psi_n-\psi|\to0.
\]
The configuration gradients of these square roots are weakly $L^2(M)$, so
their products identify the full diffusion flux
\[
 2M\sqrt{\psi_n/M}\,\grad_q\sqrt{\psi_n/M}
       \longrightarrow M\grad_q(\psi/M)
\]
against bounded smooth tests up to the closed ball. The zero-flux variational
identity therefore passes with its original boundary meaning. The spatial
diffusion term vanishes at order $\varepsilon_n$ by normalization.

Fluid convection passes by strong $L^2$ velocity convergence, viscosity by
weak $H^1$ convergence, and forcing by strong stress convergence. Recover
periodic pressure from the limiting divergence equation. The common initial
terms pass in the time-integrated weak equations. Smooth test pairings are
continuous in time by the flux bounds. Entropy uniform integrability extends
this to weak $L^1$ density continuity: approximate a bounded measurable test
by bounded smooth functions outside a set of arbitrarily small measure, then
use the common uniform integrability on that set. The velocity energy class
gives its weak $L^2$ time representative. Bounded $\log g_0$ and the initial
entropy inequality yield the strong initial density trace by the relative
entropy argument in Section~\ref{subsec:identification}.

Take an almost-everywhere-time subsequence of the strong convergences and use
convex lower semicontinuity of entropy, viscosity, and configuration Fisher
dissipation in \eqref{eq:entropy}. Drop its nonnegative spatial Fisher term.
This gives the usual limiting initial-time inequality. If the positive family
was selected to have almost-everywhere time-shifted inequalities, the higher
conditional density moment makes the entropy densities uniformly integrable,
and those inequalities pass too; that optional selection is not needed above.
Finally use nested subsequences on integer time intervals and diagonalize.
They define one global pair and one decreasing $\varepsilon_j\to0$ with all
the conclusions of Theorem~\ref{thm:main}. The fluid viscosity, FENE potential,
full velocity-gradient drag, and zero configuration-flux condition have
retained their original values and meanings throughout the limit.

\paragraph{AI use and materials.}
OpenAI Codex assisted with the argument, source checks, and mathematical
exposition. Separate fresh Codex contexts supplied an AI reader trace and an
AI referee report on the preceding draft; this revision responds to those
reports. Independent expert review is requested. The manuscript contains the
complete submitted analytic argument and uses no proof-critical numerical
computation. Its editable \LaTeX\ source accompanies the submission.

\begin{thebibliography}{9}
\bibitem{BS} J. W. Barrett and E. S\"uli,
\emph{Existence and equilibration of global weak solutions to finitely
extensible nonlinear bead-spring chain models for dilute polymers},
arXiv:1004.1432v2 (2011).
\href{https://arxiv.org/abs/1004.1432v2}{arXiv:1004.1432v2}.
\bibitem{Masmoudi} N. Masmoudi,
\emph{Global existence of weak solutions to the FENE dumbbell model of
polymeric flows}, Invent. Math. \textbf{191} (2013), 427--500.
\href{https://arxiv.org/abs/1004.4015}{arXiv:1004.4015}.
The configuration estimates and appendix proposition locators used here refer
to the revised \href{https://math.nyu.edu/~masmoudi/Fene-global.pdf}{author manuscript}.
\bibitem{BCC} P. Bonicatto, G. Ciampa, and G. Crippa,
\emph{On the advection-diffusion equation with rough coefficients:
weak solutions and vanishing viscosity}, arXiv:2107.03659.
\href{https://arxiv.org/abs/2107.03659}{arXiv:2107.03659}.
\bibitem{DS} T. D\k{e}biec and E. S\"uli,
\emph{On a class of generalised solutions to the kinetic Hookean dumbbell
model for incompressible dilute polymeric fluids: existence and macroscopic
closure}, Arch. Ration. Mech. Anal. \textbf{249} (2025), article 43.
\href{https://doi.org/10.1007/s00205-025-02115-x}{doi:10.1007/s00205-025-02115-x}.
\end{thebibliography}
\end{document}
